\documentclass[12pt]{article}
\usepackage[utf8]{inputenc}
\usepackage[T1]{fontenc}
\usepackage{hyperref}
\title{Moroccan Sign Language Language Resources for K-12 Curriculum Development}
\author{Compiled by Claude Sonnet 5.5 (claude-sonnet-5-5) for Andr\'as Kornai}
\date{}
\begin{document}
\maketitle

\section{Introduction}
Moroccan Sign Language (MSL, ISO 639-3 xms) is the sign language of the deaf community in Morocco. The search results report that schooling for deaf children is poor: one source says there is no secondary education for the deaf, and an older project page says about 85 per cent of deaf or hard-of-hearing children do not attend school. One source also says Morocco has several regional varieties rather than a single unified sign language. No official national MSL curriculum or textbook was found.

\section{Primary Educational Resources}
\begin{itemize}
\item ``Let's Communicate'' is described as Morocco's first online sign language dictionary, with videos and interactive visual aids. Its founder is a former Professional Fellows participant. Its current status is not confirmed. \url{https://www.handsalongthenile.org/mounir-launches-lets-communicate/}
\item A Zero Project entry describes a project on improving deaf children's reading through technology. \url{https://zeroproject.org/view/project/b1d89655-9917-eb11-a813-0022489b3a6d}
\end{itemize}

\section{Linguistic and Reference Materials}
\begin{itemize}
\item A conference paper in the Hamburg sign-language LREC proceedings, from the search results. \url{https://www.sign-lang.uni-hamburg.de/lrec/pub/16033.pdf}
\item Work by Abdelhadi Soudi and colleagues on an MSL STEM thesaurus, linked to an Arabic concordancer. \url{https://qscience.com/content/papers/10.5339/qfarf.2013.ICTP-05}
\end{itemize}

\section{Religious and Community Resources}
The IFES electoral lexicon was adapted into MSL with an MSL consultant. \url{https://ifes.org/news/video-ifes-creates-electoral-lexicon-moroccan-sign-language} A second Hands Along the Nile page from the results is \url{https://www.handsalongthenile.org/?p=1943}. No Moroccan deaf association site was surfaced by the search, and no scripture materials were found.

\section{Digital and Audio Resources}
A 2025 paper describes a rule-based machine-translation system that generates MSL from Arabic text. \url{https://aclanthology.org/2025.bucc-1.1.pdf} Further material appears on the National Academies project page. \url{https://sites.nationalacademies.org/pga/peer/PGA_152065}

\section{Curriculum Development Context}
Some video and dictionary tools exist, but no complete curriculum was found. A K-12 programme would need to be built with Moroccan deaf communities and teachers, and the regional variation reported in the sources would need to be addressed.
\end{document}